% Pass5388--5392: all-q Levi flag distance/Hodge theorem.
\subsection{All-\(q\) Levi flag distance scheme and Hodge cycle projector}

Let \(\mathcal S\) be a finite generalized quadrangle of order \((q,q)\), \(q>1\),
let \(\Gamma\) be its point--line incidence graph, and let
\[
X=L(\Gamma)
\]
be the line graph of \(\Gamma\).  Thus the vertices of \(X\) are exactly the flags
\((p,L)\) of \(\mathcal S\), and two flags are adjacent when they share their point or
their line.

There is an important prior-art boundary.  Colangelo, Monzillo, and Siciliano,
\emph{Discrete Mathematics} \textbf{347} (2024), 114054 (arXiv:2406.03942), already
construct and characterize the symmetric primitive four-class fusion of the flag
association scheme.  In the equal-parameter case its nontrivial valencies are
\[
2q,\qquad 2q^2,\qquad 2q^3,\qquad q^4.
\]
The result below does not claim priority for that association scheme.  The new bridge here
is to the Levi Hodge/cycle projector used throughout the \(W(3,3)\) architecture.

Direct generalized-quadrangle counting identifies the four fusion relations with the four
distance shells of \(X\), giving the intersection array
\[
\boxed{
\{2q,q,q,q;\;1,1,1,2\}.
}
\]
Hence every flag has distance profile
\[
\boxed{
1,\;2q,\;2q^2,\;2q^3,\;q^4,
}
\]
and
\[
|V(X)|=(q+1)^2(q^2+1).
\]
The corresponding distance-quotient characteristic polynomial factors exactly as
\[
\boxed{
(x-2q)(x+2)(x-q+1)\bigl((x-q+1)^2-2q\bigr).
}
\]
Thus the adjacency spectrum of \(X\) is
\[
2q,\qquad q-1+\sqrt{2q},\qquad q-1,\qquad
q-1-\sqrt{2q},\qquad -2,
\]
with multiplicities
\[
\boxed{
1,\quad \frac{q(q+1)^2}{2},\quad q(q^2+1),\quad
\frac{q(q+1)^2}{2},\quad q^4.
}
\]

The multiplicity \(q^4\) has a direct topological meaning.  The Levi graph has
\[
|V(\Gamma)|=2(q+1)(q^2+1),
\qquad
|E(\Gamma)|=(q+1)^2(q^2+1),
\]
so, because \(\Gamma\) is connected,
\[
\boxed{
\beta_1(\Gamma)=|E|-|V|+1=q^4.
}
\]
Orient every Levi edge from its point endpoint to its line endpoint and let \(D\) be the
oriented incidence matrix.  Two distinct oriented columns have inner product one precisely
when the corresponding Levi edges meet, hence
\[
\boxed{
D^{\mathsf T}D=2I+A_X.
}
\]
Consequently
\[
\boxed{
\ker D=E_{-2}(A_X),
}
\]
so the terminal primitive eigenspace of the flag distance scheme is \emph{exactly} the
Levi cycle space.  Its dimension is \(q^4\).

Let \(A_0,\ldots,A_4\) be the distance matrices of \(X\).  The standard sequence for
the eigenvalue \(-2\) is
\[
1,-\frac1q,\frac1{q^2},-\frac1{q^3},\frac1{q^4}.
\]
Multiplying by the distance valencies gives the \(q\)-independent terminal first-eigenmatrix
row
\[
\boxed{(1,-2,2,-2,1),}
\]
while multiplying by the multiplicity \(q^4\) gives the terminal second-eigenmatrix column
\[
\boxed{(q^4,-q^3,q^2,-q,1).}
\]
Therefore the orthogonal Hodge projector onto the Levi cycle space is the radial Bose--Mesner
idempotent
\[
\boxed{
P_{\rm cyc}=E_{-2}
=
\frac{q^4A_0-q^3A_1+q^2A_2-qA_3+A_4}
{(q+1)^2(q^2+1)}.
}
\]

At \(q=3\), this becomes
\[
|V(X)|=160,
\qquad
(1,6,18,54,81),
\qquad
\beta_1(\Gamma)=81,
\]
and
\[
\boxed{
E_{-2}
=
\frac{81A_0-27A_1+9A_2-3A_3+A_4}{160}.
}
\]
This is exactly the BT545--BT551 W33 Levi Hodge/Kirchhoff projector and explains the
previously observed coefficient column \((81,-27,9,-3,1)\) as the \(q=3\) specialization
of the universal \((-1/q)^d\) radial law.

\paragraph{Evidence boundary.}
This theorem identifies an exact graph-theoretic/Hodge sector.  It does not by itself prove
a quantum-code distance, a fault-tolerance threshold, or a laboratory protection mechanism.
